\section*{Chapter \ref{ch:iv:estimation} --- Estimation}

\begin{solution}{iq:iv:estimation:1}
400 people $\times$ 2 cups a day $\times$ 220 working days $\times$ 0.25 litre $= 44\,000$ litres a year; a 90\%
interval of a factor of about 2 either way (15\,000 to 90\,000), since the cups per person is the weakest
factor. There is no published figure; the question scores the chain and the units.
\iqlookfor{a clean chain with units and a stated interval.}
\end{solution}

\begin{solution}{iq:iv:estimation:2}
$365 \times 24 \times 3\,600 = 31\,536\,000$, about $3.15 \times 10^7$ (``$\pi \times 10^7$'' is a good
mnemonic). Regular US equity trading: $252 \times 6.5 \times 3\,600 = 5\,896\,800$, about $5.9 \times 10^6$,
under a fifth of the year.
\iqlookfor{exact arithmetic with the right calendar and session length.}
\end{solution}

\begin{solution}{iq:iv:estimation:3}
2\,000, the geometric mean. Bounds a factor of 100 apart describe uncertainty on the log scale; the midpoint of
the logs is the geometric mean, and it is a factor of 10 from each bound. The arithmetic mean 10\,100 is only a
factor of 2 from the upper bound and 50 from the lower, which treats a factor-of-50 miss on one side like a
factor-of-2 miss on the other.
\iqlookfor{thinking in logs for multiplicative uncertainty.}
\end{solution}

\begin{solution}{iq:iv:estimation:4}
A chain, with each factor an assumption to be checked: suppose world trade in goods and services is of the
order of 30 trillion dollars a year; FX turnover is dominated by
financial flows and rolling of swaps, perhaps 50 to 100 times the trade flows a year, over 250 days: 6 to 12
trillion a day. A 90\% interval of 3 to 15 trillion (geometric centre 6.7). The BIS survey measured 9.6 trillion
a day in April 2025 (\cref{dat:iv:estimation:anchors}), inside the interval.
\iqlookfor{a chain from a known macro quantity, a multiplier for financial activity, and a wide honest interval.}
\end{solution}

\begin{solution}{iq:iv:estimation:5}
The chain of \cref{ex:iv:estimation:opra} gives about 850\,000; with each factor uncertain by a factor of 2, a
90\% interval of 200\,000 to 2 million. The July 2026 peak was 0.89 million (\cref{dat:iv:estimation:anchors}).
Before building: ask for the published peak statistics and capacity projections, the message format (bytes per
message, packets per message), the number of lines and their split, and the growth rate, because systems are
sized for next year's burst.
\iqlookfor{a structural chain, a wide interval, and the engineer's instinct to ask for the published capacity.}
\end{solution}

\begin{solution}{iq:iv:estimation:6}
Perhaps 50 actively traded products (interest-rate futures and options, equity index, energy, metals,
agriculture) at 500\,000 contracts a day on average, with a long tail: about 25 million, 90\% interval 10 to 60
million. The group's exchanges averaged 28.1 million contracts a day in 2025.
\iqlookfor{a count-times-activity chain and an interval that contains the answer without being useless.}
\end{solution}

\begin{solution}{iq:iv:estimation:7}
About 8.2 billion people divided by 365.25 days: about 22 million. Birthdays are not quite uniform over the year
(seasonal birth patterns), which moves the answer by some percent, not a factor.
\iqlookfor{the division, and awareness of the uniformity assumption.}
\end{solution}

\begin{solution}{iq:iv:estimation:8}
About 5\,000 stocks $\times$ 20\,000 top-of-book changes a day on average (many more for the most active, far
fewer for the tail) $\times$ 32 bytes (timestamp, symbol identifier, bid and ask prices and sizes) $= 3.2
\times 10^9$ bytes, about 3 GB a day. The distribution across stocks is very skewed, so the per-stock average
is the weak factor; a second chain from total message counts would check it.
\iqlookfor{a record-size estimate and awareness of skew across symbols.}
\end{solution}

\begin{solution}{iq:iv:estimation:9}
If each interval contains the truth with probability 0.9 independently, the number of hits is binomial(10,
0.9), and $\P(X \le 5) \approx 0.0016$: very surprising. The intervals are too narrow. Widen them, roughly
doubling their log width if the misses were large (\cref{fig:iv:estimation:calib}), and keep a record of hits
to recalibrate.
\iqlookfor{a binomial test of calibration and the practical correction.}
\end{solution}

\begin{solution}{iq:iv:estimation:10}
A factor of 2 at 90\% means a log standard deviation of $\ln 2/1.645 \approx 0.42$ per factor. Four independent
factors give $0.42 \times 2 = 0.84$, so the product is within a factor of $e^{1.645 \times 0.84} = 4$ of the truth
at 90\%, far less than $2^4 = 16$, because the four errors are unlikely to go the same way at once.
\iqlookfor{adding log variances, and why independent errors partly cancel.}
\end{solution}

\begin{solution}{iq:iv:estimation:11}
Weights $1/0.7^2 \approx 2.04$ and $1/1.0^2 = 1$ on the logs give $\exp\big((2.04 \ln 10^6 + \ln 10^7)/3.04\big)
\approx 2.1 \times 10^6$, with log standard deviation $1/\sqrt{3.04} \approx 0.57$. Refuse to combine when the
chains disagree by more than their uncertainties allow (here the difference of the logs, 2.3, is about 1.9
standard deviations of the difference, $\sqrt{0.49 + 1} \approx 1.22$: borderline) or when they share a factor, so
that their errors are not independent. Then look for the wrong factor.
\iqlookfor{inverse-variance weighting on the log scale, and a test of consistency before combining.}
\end{solution}

\begin{solution}{iq:iv:estimation:12}
With $\alpha = 0.1$: $[5, 12]$ contains 9.6 and scores its width, 7. $[2, 5]$ misses by 4.6 and scores $3 + 20
\times 4.6 = 95$. The rule is proper: for any belief about the quantity, the expected score is minimised by
reporting the belief's own 5\% and 95\% quantiles, so shading the interval narrower (or wider) than one believes
raises the expected penalty (Gneiting and Raftery, 2007).
\iqlookfor{computing the score and knowing why a proper rule makes honesty optimal.}
\end{solution}

\begin{solution}{iq:iv:estimation:13}
$63.9 \times 10^6 \times 40 \times 8 \approx 20.4$ gigabits a second. The capacity projection is 4.403 gigabits
per 100 milliseconds, 44 gigabits a second: about twice the estimate. The estimate leaves out packet and
protocol overhead, the fact that a one-second peak averages over shorter and higher bursts (the 10 ms peak of
0.89 million messages is 89 million a second), retransmission and a second, redundant stream, and growth.
\iqlookfor{a units-careful bandwidth calculation and the reasons capacity exceeds the average peak.}
\end{solution}
